\documentclass[10pt,a4paper]{amsart}
\usepackage[margin=19mm]{geometry}
\usepackage{amsmath,amssymb,mathtools}
\usepackage[scr=rsfs,cal=euler]{mathalfa}
\usepackage{xfrac}
\usepackage[unicode,hidelinks,bookmarks=false]{hyperref}
\newcommand{\ie}{i.e.\ }
\newcommand{\cf}{cf.\ }
\newcommand{\resp}{respectively\ }
\newcommand{\Subsection}{\subsection}
\providecommand{\nobreakdash}{\nobreak\hbox{-}}
\setlength{\emergencystretch}{2em}
\multlinegap0pt
\usepackage{amssymb}
\usepackage{enumerate}
\newtheorem{lemma}{Lemma}
\title{Analytic regularity for a fourth-order singularly perturbed boundary balue problem with two small parameters}
\author{I. Sykopetritou, C. Xenophontos}
\date{Source: arXiv:2607.16781v1 (CC BY 4.0)}
\begin{document}
\maketitle

\noindent\emph{Source and licence.} Analytic regularity for a fourth-order singularly perturbed boundary balue problem with two small parameters. Version arXiv:2607.16781v1 (CC BY 4.0), distributed by arXiv under the Creative Commons Attribution 4.0 International licence, http://creativecommons.org/licenses/by/4.0/, as recorded in the arXiv licence metadata for this exact version and shown on the abstract page https://arxiv.org/abs/2607.16781v1. Full licence notice: \texttt{licenses/arxiv-2607.16781-CC-BY-4.0.txt}.\par\smallskip

\noindent\textit{Editorial scope.} Counted unit: the article's lemma lem:ut\_ij (source0.tex lines 538--549). The article's complete original proof of it is delivered with it (source0.tex lines 551--718, one \texttt{proof} environment of 168 lines). No mathematics is written, repaired or re-derived: the statement and the proof are the article's own lines, in the article's own notation, and no retained line is substituted or re-flowed.  Retained uncounted as the source-local context the proof needs: lines 135--140; lines 148--150; lines 284--299; lines 312--377; lines 421--429; lines 434--464.  Nothing was omitted from any retained span, so the package carries no omission record.  No retained line contains a citation, so the wrapper carries no bibliography and no bibliography entry.  No retained line contains a figure environment, an image-inclusion command or drawing code, so no image is delivered and the package declares no asset dependency.  Editorial formatting only, in addition: an AMS \texttt{amsart} preamble that restates the article's own declarations and macros used by the retained lines, namely the environments \texttt{lemma} on separate counters, together with the standard packages those lines need.  The retained environments print in this document as Lemma 1 in order of appearance, whereas the article prints the same items with the numbering of its own full document.  The complete native source of the article as distributed in the arXiv e-print for this exact version is bundled unchanged as \texttt{sources/205498/source0.tex}.
\begin{equation}
\left\Vert f^{(n)}\right\Vert _{\infty ,I}\lesssim n!\gamma
_{f}^{n},\left\Vert c^{(n)}\right\Vert _{\infty ,I}\lesssim n!\gamma
_{c}^{n},\left\Vert b^{(n)}\right\Vert _{\infty ,I}\lesssim n!\gamma
_{b}^{n}\;\forall \;n\in \mathbb{N}_{0}.  \label{analytic}
\end{equation}%

\begin{equation}
\varepsilon _{1} \ll \varepsilon _{2}^{2},  \label{eps_relation}
\end{equation}%

For $u_{i,j}$, i.e. for the (slow) variable $x$, we obtain
%
\begin{equation}
\left. 
\begin{array}{c}
u_{0,0} = \frac{f}{c} \\
u_{1,0}=0,u_{i,0}=\frac{\left( bu_{i-2,0}^{\prime }\right) ^{\prime }}{c}%
,i\geq 2 \\  
u_{i,1}=0,i\geq 0 \\
u_{0,j}=u_{1,j}=u_{2,j}=u_{3,j}=0,j\geq 2 \\
u_{i,j}=\frac{1}{c}\left\{ \left( bu_{i-2,j}^{\prime }\right) ^{\prime
}-u_{i-4,j-2}^{(4)}\right\} ,i\geq 4,j\geq 2%
\end{array}%
\right\},  \label{uij}
\end{equation}
with the details appearing in the Appendix.

\begin{equation}
\left. 
\begin{array}{c}
-\tilde{b}_{0}\left(\tilde{u}^{BL}_{0,0}\right)^{\prime \prime }+\tilde{c}_{0}\tilde{u}^{BL}_{0,0}=0 \\
-\tilde{b}_{0}\left(\tilde{u}^{BL}_{i,0}\right)^{\prime \prime }+\tilde{c}_{0}\tilde{u}^{BL}_{i,0}=%
\sum_{\ell =1}^{i}\left\{ \left( \tilde{b}_{\ell }%
\left(\tilde{u}^{BL}_{i-\ell ,0}\right)^{\prime }\right) ^{\prime }-\tilde{c}_{\ell }\tilde{u}^{BL}_{i-\ell ,0}\right\} ,i\geq 1 \\
-\tilde{b}_{0}\left(\tilde{u}^{BL}_{0,1}\right)^{\prime \prime }+\tilde{c}_{0}\tilde{u}^{BL}_{0,1}=0 \\
-\tilde{b}_{0}\left(\tilde{u}^{BL}_{i,1}\right)^{\prime \prime }+\tilde{c}_{0}\tilde{u}^{BL}_{i,1}=%
\sum_{\ell =1}^{i}\left\{ \left( \tilde{b}_{\ell }\left(\tilde{u}^{BL}_{i-\ell ,1}\right)^{\prime }\right) ^{\prime }-\tilde{c}_{\ell }\tilde{u}^{BL}_{i-\ell ,1}\right\} ,i\geq 1 \\
-\tilde{b}_{0}\left(\tilde{u}^{BL}_{0,j}\right)^{\prime \prime }+\tilde{c}_{0}\tilde{u}^{BL}_{0,j}=-\left(\tilde{u}^{BL}_{0,j-2}\right)^{(4)},j\geq 2 \\
-\tilde{b}_{0}\left(\tilde{u}^{BL}_{i,j}\right)^{\prime \prime }+\tilde{c}_{0}\tilde{u}^{BL}_{i,j}=-\left(\tilde{u}^{BL}_{i,j-2}\right)^{(4)}+\sum_{\ell =1}^{i}\left\{ \left( \tilde{b}_{\ell }%
\left(\tilde{u}^{BL}_{i-\ell ,j}\right)^{\prime }\right) ^{\prime }-\tilde{c}_{\ell }\tilde{u}^{BL}_{i-\ell ,j}\right\} ,i\geq 1,j\geq 2%
\end{array}%
\right\}.  \label{ut_ij}
\end{equation}
%
Correspondingly, for the boundary layer $\bar{u}_{i,j}^{BL}$, i.e. for the (fast)
variable $\bar{x}=x\varepsilon _{2}/\varepsilon _{1}$, we first write 
$b(x)=b(\varepsilon _{1}\varepsilon _{2}x/\varepsilon
_{1}\varepsilon _{2})=b(\varepsilon _{1}\bar{x}/\varepsilon _{2})$ and use a Taylor
series to get
%
\begin{equation*}
b(\bar{x})=\sum_{\ell =0}^{\infty } \left( \frac{\varepsilon _{1}}{\varepsilon _{2}}\right) ^{\ell }\bar{b}_{\ell }(\bar{x}),
\end{equation*}
%
where $\bar{b}_{\ell }(\bar{x})=$ $\frac{b^{(\ell )}(0)}{\ell !}\bar{x}^{\ell }$, and
similarly for $c$. Then, we find that (see the Appendix for the details)
%
\begin{equation}
\left. 
\begin{array}{c}
\left(\bar{u}^{BL}_{i,0}\right)^{(4)}-\bar{b}_{0}\left(\bar{u}^{BL}_{i,0}\right)^{\prime \prime }=0,i\geq 0 \\
\left(\bar{u}^{BL}_{0,1}\right)^{(4)}-\bar{b}_{0}\left(\bar{u}^{BL}_{0,1}\right)^{\prime \prime }=0 \\
\left(\bar{u}^{BL}_{i,1}\right)^{(4)}-\bar{b}_{0}\left(\bar{u}^{BL}_{i,1}\right)^{\prime \prime }=\left( \bar{b}_{1}\left(\bar{u}^{BL}_{i-1,0}\right)^{\prime }\right) ^{\prime },i\geq 1 \\
\left(\bar{u}^{BL}_{0,j}\right)^{(4)}-\bar{b}_{0}\left(\bar{u}^{BL}_{0,j}\right)^{\prime \prime }=-\bar{c}_{0}\bar{u}^{BL}_{0,j-2},j\geq 2 \\
\left(\bar{u}^{BL}_{1,2}\right)^{(4)}-\bar{b}_{0}\left(\bar{u}^{BL}_{1,2}\right)^{\prime \prime }=-\bar{c}_{0}\bar{u}^{BL}_{1,0}+\left( \bar{b}_{1}\left(\bar{u}^{BL}_{0,1}\right)^{\prime }\right) ^{\prime } \\
\left(\bar{u}^{BL}_{i,2}\right)^{(4)}-\bar{b}_{0}\left(\bar{u}^{BL}_{i,2}\right)^{\prime \prime }=-\bar{c}_{0}\bar{u}^{BL}_{i,0}+\left( \bar{b}_{2}\left(\bar{u}^{BL}_{i-2,0}\right)^{\prime }\right) ^{\prime }+\left( \bar{b}_{1}\left(\bar{u}^{BL}_{i-1,1}\right)^{\prime }\right) ^{\prime },i\geq 2 \\
\left(\bar{u}^{BL}_{i,j}\right)^{(4)}-\bar{b}_{0}\left(\bar{u}^{BL}_{i,j}\right)^{\prime \prime }=-\bar{c}_{0}\bar{u}^{BL}_{i,j-2}+\left( \bar{b}_{j}\left(\bar{u}^{BL}_{i-j,0}\right)^{\prime }\right) ^{\prime }+\left( \bar{b}_{j-1}\left(\bar{u}^{BL}_{i-j+1,1}\right)^{\prime }\right) ^{\prime }+ \\
+\sum_{\ell =1}^{j-2}\left\{ \left( \bar{b}_{\ell}\left(\bar{u}^{BL}_{i-\ell ,j-\ell}\right)^{\prime }\right) ^{\prime }-\bar{c}_{\ell}\bar{u}^{BL}_{i-\ell ,j-\ell -2}\right\} ,i\geq 3,j=3,\ldots ,i \\
\left(\bar{u}^{BL}_{i,j}\right)^{(4)}-\bar{b}_{0}\left(\bar{u}^{BL}_{i,j}\right)^{\prime \prime }=-\bar{c}_{0}\bar{u}^{BL}_{i,j-2}+\left( \bar{b}_{i}\left(\bar{u}^{BL}_{0,1}\right)^{\prime }\right) ^{\prime }+ \\
+\sum_{\ell =1}^{i-1}\left\{ \left( \bar{b}_{\ell}\left(\bar{u}^{BL}_{i-\ell ,j-\ell}\right)^{\prime }\right) ^{\prime }-\bar{c}_{\ell}\bar{u}^{BL}_{i-\ell ,j-\ell -2}\right\} ,i\geq 2,j=i+1 \\
\left(\bar{u}^{BL}_{i,j}\right)^{(4)}-\bar{b}_{0}\left(\bar{u}^{BL}_{i,j}\right)^{\prime \prime }=-\bar{c}_{0}\bar{u}^{BL}_{i,j-2}+ \\
+\sum_{\ell =1}^{i}\left\{ \left( \bar{b}_{\ell}\left(\bar{u}^{BL}_{i-\ell ,j-\ell}\right)^{\prime }\right) ^{\prime }-\bar{c}_{\ell}\bar{u}^{BL}_{i-\ell ,j-\ell -2}\right\} ,i\geq 1,j>i+1
\end{array}%
\right\}.  \label{ubar_ij}
\end{equation}
%
Both (\ref{ut_ij}) and (\ref{ubar_ij}) are supplemented with the boundary
conditions$\;\forall \;i,j:$
%
\begin{equation}
\left. 
\begin{array}{c}
\tilde{u}^{BL}_{i,j}(0)+\bar{u}^{BL}_{i,j}(0)=-u_{i,j}(0) \\ 
\left(\tilde{u}^{BL}_{i,j}\right)^{\prime }(0)+\left(\bar{u}^{BL}_{i,j}\right)^{\prime }(0)=-u_{i,j}^{\prime }(0)
\\ 
\lim_{x\rightarrow \infty }\tilde{u}^{BL}_{i,j}(x)=0 \; , \; \lim_{x\rightarrow \infty }\bar{u}^{BL}_{i,j}(x)=0 \\
\lim_{x\rightarrow \infty }\left(\tilde{u}^{BL}_{i,j}\right)^{\prime }(x)=0 \; , \; \lim_{x\rightarrow \infty }\left(\bar{u}^{BL}_{i,j}\right)^{\prime }(x)=0
\end{array}
%
\right\} .  \label{bcBL}
\end{equation}
%
Analogous problems are satisfied by $\check{u}^{BL}_{i,j}$ and $\hat{u}^{BL}_{i,j}$. A remark is in order regarding the boundary conditions involving limits: they are meant to enforce \emph{decay} in the functions, even though strictly speaking, they are not defined on $(0, \infty)$.

\begin{lemma}\label{lem:u_n_ij}
Let $u_{i,j}$ be defined by (\ref{uij}) and assume  (\ref{analytic})--(\ref{eps_relation}) hold. Then there exist
positive constants $C, K_1, K_2,$ such that
%
\begin{equation}
\left\|u_{i, j}^{(n)}\right\|_{\infty, I} \leq C n ! K_1^n i^i K_2^i \quad \forall n \in \mathbb{N} .
\end{equation}
%
\end{lemma}

\begin{lemma}
\label{lemma_aux0} Let $\lambda_1 \in \mathbb{C}$ with $Re(\lambda_1)>0,Re(\lambda_1^{2})>0$. Let $F_1$ be an entire function satisfy- ing, for some $%
C_{F_1}>0,$ $j\in \mathbb{N}_{0},$ $q_1\geq (j+1/2)/|\lambda_1|>0,$ 
\begin{equation*}
\left\vert F_1(z)\right\vert \leq C_{F_1}e^{-Re(\lambda_1 z)}\left( q_1+\left\vert
z\right\vert \right) ^{j} \; \; \forall \; z\in \mathbb{C}.
\end{equation*}%
Similarly, let $j\in \mathbb{N}, \kappa, \lambda_2 \in \mathbb{C}^{+},$ with $\lambda_2 = \kappa^2$ and let $F_2$ be an entire function satisfying, for some $C_{F_2}>0,$ $q_2\geq \frac{4j}{|\kappa|}>0,$
\begin{equation*}
\left\vert F_2(z)\right\vert \leq C_{F_2}e^{-Re(\kappa z)}\left( q_2+\left\vert
z\right\vert \right) ^{2j-1} \; \; \forall \; z\in \mathbb{C}.
\end{equation*}%
With $\alpha_1, \alpha_2 \in \mathbb{C}$, let $v,w:(0,\infty )\rightarrow \mathbb{C}$ be the solutions of the problems
%
\begin{equation*}
-v^{\prime \prime }+\lambda_1^{2}v=F_1\text{ on }(0,\infty
) \; , \; v(0)=\alpha_1 - w(0) \; , \; \lim_{x\rightarrow \infty }v(x)=0,
\end{equation*}%
\begin{equation*}
w^{(4)}-\lambda_2 w^{(2)}=F_2\text{ on }(0,\infty ) \; , \; w^{\prime }(0)=\alpha_2 - v^{\prime }(0) \; , \; \lim_{x\rightarrow \infty}w(x)=0.
\end{equation*}%
Then $v,w$ can be extended to entire functions (denoted again by $v,w$) which satisfy
\begin{equation*}
\left\vert v(z)\right\vert \leq C \left[ \frac{C_{F}}{j+1}\left(
q_*+\left\vert z\right\vert \right) ^{j+1}+\left\vert \alpha_1 \right\vert+\left\vert \alpha_2 \right\vert \right] e^{-Re(\lambda_1 z)} \; \; \forall \; z\in \mathbb{C},
\end{equation*}%
\begin{equation*}
\left\vert w(z)\right\vert \leq C \left[ \frac{C_{F}}{2j}\left( q_*+\left\vert z\right\vert \right) ^{2j}+\left\vert \alpha_1 \right\vert+\left\vert \alpha_2 \right\vert \right] e^{-Re(\kappa z)} \; \; \forall \; z\in \mathbb{C},
\end{equation*}%
where $C>0,q_* = \max \{q_1,q_2^2 \},C_F = \max \{C_{F_1},C_{F_2} \}$.
\end{lemma}

\begin{lemma}
\label{lem:ut_ij}
The functions $\tilde{u}^{BL}_{i,j}, \bar{u}^{BL}_{i,j}$ which satisfy (\ref{ut_ij})--(\ref{bcBL}), are entire and there exist
positive constants $C, \tilde{\gamma}, \bar{\gamma},$ such that
%
\begin{eqnarray*}
\vert \tilde{u}^{BL}_{i,j}(z) \vert &\leq& C \tilde{\gamma}^{i+j} \frac{(2i+j+|z|)^{2i+j}}{i!} e^{-\tilde{\beta} Re(z)} \; \forall \; z \in \mathbb{C}, \; Re(z) > 0,\\
\vert \bar{u}^{BL}_{i,j}(z) \vert &\leq& C \bar{\gamma}^{i+j} \frac{(i+2j+|z|)^{i+2j}}{j!} e^{-\bar{\beta} Re(z)} \; \forall \; z \in \mathbb{C}, \; Re(z) > 0,
\end{eqnarray*}
%
where $\tilde{\beta}=\sqrt{\tilde{c}_{0} / \tilde{b}_{0}}$ and $\bar{\beta}=\sqrt{\bar{b}_{0}}$.
\end{lemma}

\begin{proof}
We recall that
\begin{equation*}
\tilde{b}_{\ell}(\tilde{x}) = \frac{\tilde{x}^{\ell} b^{(\ell)}(0)}{\ell !} \; , \; \tilde{c}_{\ell}(\tilde{x}) = \frac{\tilde{x}^{\ell} c^{(\ell)}(0)}{\ell !} \; , \; \bar{b}_{\ell}(\bar{x}) = \frac{\bar{x}^{\ell} b^{(\ell)}(0)}{\ell !} \; , \; \bar{c}_{\ell}(\bar{x}) = \frac{\bar{x}^{\ell} c^{(\ell)}(0)}{\ell !}.
\end{equation*}
Consequently, there exist positive constants $C_{\tilde{b}}, \gamma _{\tilde{b}}, C_{\tilde{c}}, \gamma _{\tilde{c}}, C_{\bar{b}}, \gamma _{\bar{b}}, C_{\bar{c}}, \gamma _{\bar{c}},$ depending only on $b,c,$ such that
\begin{equation*}
|\tilde{b}_{\ell}(z)| \leq C_{\tilde{b}} \gamma _{\tilde{b}} |z|^{\ell} \; , \; |\tilde{c}_{\ell}(z)| \leq C_{\tilde{c}} \gamma _{\tilde{c}} |z|^{\ell} \; , \; |\bar{b}_{\ell}(z)| \leq C_{\bar{b}} \gamma _{\bar{b}} |z|^{\ell} \; , \; |\bar{c}_{\ell}(z)| \leq C_{\bar{c}} \gamma _{\bar{c}} |z|^{\ell}.
\end{equation*}
Then, with $K_{2}$ the constant from Lemma \ref{lem:u_n_ij}, we choose $\tilde{\gamma } > \max \{ K_{2}, \gamma _{\tilde{b}}, \gamma _{\tilde{c}} \}$,$\bar{\gamma } > \max \{ K_{2}, \gamma _{\bar{b}}, \gamma _{\bar{c}} \}$ so that
\begin{equation*}
\left[ \frac{\gamma _{\tilde{b}} / \tilde{\gamma }}{1 - \gamma _{\tilde{b}} / \tilde{\gamma }} + \frac{\gamma _{\tilde{c}} / \tilde{\gamma }}{1 - \gamma _{\tilde{c}} / \tilde{\gamma }} \right] < 1
\end{equation*}
and
\begin{equation*}
\left[ \frac{\gamma _{\bar{b}} / \bar{\gamma }}{1 - \gamma _{\bar{b}} / \bar{\gamma }} + \frac{\gamma _{\bar{c}} / \bar{\gamma }}{1 - \gamma _{\bar{c}} / \bar{\gamma }} \right] < 1.
\end{equation*}
%
From (\ref{ut_ij}) and (\ref{ubar_ij}), we have that $\tilde{u}^{BL}_{0,0}(\tilde{x}),\bar{u}^{BL}_{0,0}(\bar{x})$ are given by 
\begin{equation*}
\tilde{u}^{BL}_{0,0}(\tilde{x})=\tilde{C}_{0,0}e^{-\sqrt{\frac{\tilde{c}_{0}}{%
\tilde{b}_{0}}}\tilde{x}}\; , \; \bar{u}^{BL}_{0,0}(\bar{x})=\bar{C}_{0,0}e^{-\sqrt{\bar{%
b}_{0}}\bar{x}}
\end{equation*}
and from (\ref{bcBL}) we find
%
\[
\tilde{C}_{0,0} = \frac{u_{0,0}(0) \sqrt{\bar{b}_0}+u'_{0,0}(0)}{\sqrt{\frac{\tilde{c}_0}{\tilde{b}_0}}-\sqrt{\bar{b}_0}} \; , \;
\bar{C}_{0,0} = \frac{-u_{0,0}(0) \sqrt{\frac{\tilde{c}_0}{\tilde{b}_0}}-u'_{0,0}(0)}{\sqrt{\frac{\tilde{c}_0}{\tilde{b}_0}}-\sqrt{\bar{b}_0}}.
\]
A similar result holds for $\tilde{u}^{BL}_{0,1}(\tilde{x})$ and $\bar{u}^{BL}_{0,1}(\bar{x})$, involving 
$\tilde{C}_{0,1}, \bar{C}_{0,1}$ given in terms of $u_{0,1}(0)$ and $u'_{0,1}(0)$, hence the desired estimates 
hold for $i=0,j=0,1$. Let us consider $i=0,j=2$: 
%
\[
-\tilde{b}_0 \left(\tilde{u}^{BL}_{0,2}\right)^{\prime \prime} + \tilde{c}_0 \tilde{u}^{BL}_{0,2} = - \left(\tilde{u}^{BL}_{0,0}\right)^{(4)}
\]
%
\[
\left( \bar{u}^{BL}_{0,2}\right)^{(4)} - \bar{b}_0 \left(\bar{u}^{BL}_{0,2} \right)^{\prime \prime}= - \bar{c}_0 \bar{u}^{BL}_{0,0}.
\]
Solving the above equations, we find
%
\[
\tilde{u}^{BL}_{0,2}(\tilde{x}) = \tilde{C}_{0,2} e^{-\sqrt{\frac{\tilde{c}_0}{\tilde{b}_0}} \tilde{x}} - A \left(\tilde{u}^{BL}_{0,0}\right)^{(4)}(\tilde{x})\; , \;
\bar{u}^{BL}_{0,2}(\bar{x}) = \bar{C}_{0,2} e^{-\sqrt{\bar{b}_0} \bar{x}} - B \bar{c}_0 \bar{u}^{BL}_{0,0}(\bar{x}),
\]
with $A, B \in \mathbb{R}$ depending only on $\tilde{b}_0, \tilde{c}_0, \tilde{C}_{0,0}, \bar{b}_0, \bar{c}_0, \bar{C}_{0,0}$. Since 
%
\[
\left(\tilde{u}^{BL}_{0,2}\right)^{(k)}(0) + \left(\bar{u}^{BL}_{0,2}\right)^{(k)}(0) = - u^{(k)}_{0,2}(0) \; , \; k=0,1,
\]
%
we can calculate the constants $\tilde{C}_{0,2}, \bar{C}_{0,2}$, in terms of $\tilde{b}_0, \tilde{c}_0, \bar{b}_0, u_{0,2}(0), u'_{0,2}(0)$. This shows that the desired results hold for $i=0,j=2$.

We proceed with 
induction on $j>2$, while keeping $i$ fixed at 0. Assuming the result holds for $j$, we will establish it
for $j+1$. The function $\tilde{u}^{BL}_{0,j+1}$ satisfies (see eq. (\ref{ut_ij}))
%
\[
-\tilde{b}_{0}\left(\tilde{u}^{BL}_{0,j+1}\right)^{\prime \prime }+\tilde{c}_{0}\tilde{u}^{BL}_{0,j+1}=-%
\left(\tilde{u}^{BL}_{0,j-1}\right)^{(4)} \; , \; j\geq 2.
\]
From the induction hypothesis and Cauchy's Integral Theorem for Derivatives, we get
%
\[
\vert \left(\tilde{u}^{BL}_{0,j-1}\right)^{(4)}(z) \vert \leq C \tilde{\gamma}^{j-1} (j-1+|z|)^{j-1} e^{-\tilde{\beta} Re(z)}
\]
and by Lemma \ref{lemma_aux0},
%
\begin{eqnarray*}
\vert \tilde{u}^{BL}_{0,j+1} (z) \vert &\leq& C e^{-\tilde{\beta} Re(z)} \left[ \tilde{\gamma}^{j-1} \frac{(j-1+|z|)^{j}}{j} +
|u_{0,j+1}(0)| + |u'_{0,j+1}(0)| \right]  \notag \\
&\leq& C e^{-\tilde{\beta} Re(z)} \tilde{\gamma}^{j+1} (j+1+|z|)^{j+1}\left[ \frac{1}{\tilde{\gamma}^2 j (j+1+|z|)} + \frac{\tilde{C}}{\tilde{\gamma}^{j+1} (j+1+|z|)^{j+1}} \right] \\
&\leq& C e^{-\tilde{\beta} Re(z)} \tilde{\gamma}^{j+1} (j+1+|z|)^{j+1},
\end{eqnarray*}
%
which gives the result for $\tilde{u}_{0,j+1}$.

%
We next show the result for $\bar{u}^{BL}_{0,j+1}$, which satisfies (see eq. (\ref{ubar_ij}))
\[
\left(\bar{u}^{BL}_{0,j+1}\right)^{(4)}-\bar{b}_{0}\left(\bar{u}^{BL}_{0,j+1}\right)^{\prime \prime }=-\bar{c}_{0}\bar{u}^{BL}_{0,j-1} \; , \; j\geq 2.
\]
%
By the induction hypothesis, we have
%
\[
\vert \bar{u}^{BL}_{0,j-1} (z) \vert \leq C \bar{\gamma }^{j-1} \frac{(2(j-1)+|z|)^{2(j-1)}}{(j-1)!} e^{-\bar{\beta} Re(z)}
\]
and using Lemma \ref{lemma_aux0}, we obtain
%
\begin{eqnarray*}
\vert \bar{u}^{BL}_{0,j+1} (z) \vert &\leq& C e^{-\bar{\beta} Re(z)} \left[ \bar{\gamma}^{j-1} \frac{(2(j-1)+|z|)^{2j-1}}{(2j-1) (j-1)!} +
|u_{0,j+1}(0)| + |u'_{0,j+1}(0)| \right]  \notag \\
&\leq& C e^{-\bar{\beta} Re(z)} \bar{\gamma}^{j+1} \frac{(2(j+1)+|z|)^{2(j+1)}}{(j+1)!} \left[ \frac{j(j+1)}{\bar{\gamma}^2 (2j-1) (2(j+1)+|z|)^3} + \right. \\
&+& \left. \frac{\bar{C} (j+1)^{j+1}}{\bar{\gamma}^{j+1} (2(j+1)+|z|)^{2(j+1)}} \right] \\
&\leq& C e^{-\bar{\beta} Re(z)} \bar{\gamma}^{j+1} \frac{(2(j+1)+|z|)^{2(j+1)}}{(j+1)!}.
\end{eqnarray*}
%

We finally consider the case $i\geq 1,j\geq 0$. We first assume that the result holds for $\tilde{u}^{BL}_{i,j}, j\geq 2$ and we will establish it for $\tilde{u}^{BL}_{i,j+1}$, with $i\geq 1$ fixed (but arbitrary). By (\ref{ut_ij}), the function $\tilde{u}^{BL}_{i,j+1}$ satisfies
%
$$
-\tilde{b}_0\left(\tilde{u}_{i, j+1}^{B L}\right)^{\prime \prime}+\tilde{c}_0 \tilde{u}_{i, j+1}^{B L}=-\left(\tilde{u}_{i, j-1}^{B L}\right)^{(4)}+\sum_{\ell=1}^i\left\{\tilde{b}_{\ell}\left(\tilde{u}_{i-\ell, j+1}^{B L}\right)^{\prime \prime}+\tilde{b}_{\ell}^{\prime}\left(\tilde{u}_{i-\ell, j+1}^{B L}\right)^{\prime}-\tilde{c}_{\ell} \tilde{u}_{i-\ell, j+1}^{B L}\right\} .
$$
%
In order to use Lemma \ref{lemma_aux0}, we need to bound the right hand side above:
%
$$
\tilde{G}(z):=-\left(\tilde{u}_{i, j-1}^{B L}\right)^{(4)}+\sum_{\ell=1}^i\left\{\tilde{b}_{\ell}\left(\tilde{u}_{i-\ell, j+1}^{B L}\right)^{\prime \prime}+\tilde{b}_{\ell}^{\prime}\left(\tilde{u}_{i-\ell, j+1}^{B L}\right)^{\prime}-\tilde{c}_{\ell} \tilde{u}_{i-\ell, j+1}^{B L}\right\} .
$$
%
By the induction hypothesis and Cauchy's Integral Theorem for Derivatives, we find
%
$$
\begin{aligned}
|\tilde{G}(z)| \leq & \left|\left(\tilde{u}_{i, j-1}^{B L}\right)^{(4)}\right|+C \sum_{\ell=1}^i\left\{\left|\tilde{b}_{\ell}\right|\left|\left(\tilde{u}_{i-\ell, j+1}^{B L}\right)^{\prime \prime}\right|+\left|\tilde{b}_{\ell}^{\prime}\right|\left|\left(\tilde{u}_{i-\ell, j+1}^{B L}\right)^{\prime}\right|+\left|\tilde{c}_{\ell}\right|\left|\tilde{u}_{i-\ell, j+1}^{B L}\right|\right\} \\
\leq & C \tilde{\gamma}^{i+j-1} \frac{(2 i+j-1+|z|)^{2 i+j-1}}{i!} e^{-\tilde{\beta} R e(z)}+C \sum_{\ell=1}^i|z|^{\ell}\left(\gamma_{\tilde{b}}^{\ell}+\gamma_{\tilde{c}}^{\ell}\right)\left|\tilde{u}_{i-\ell, j+1}^{B L}\right| \\
\leq & C \tilde{\gamma}^{i+j-1} \frac{(2 i+j-1+|z|)^{2 i+j-1}}{i!} e^{-\tilde{\beta} R e(z)}+ \\
& +C \sum_{\ell=1}^i|z|^{\ell}\left(\gamma_{\bar{b}}^{\ell}+\gamma_{\tilde{c}}^{\ell}\right) \tilde{\gamma}^{i+j+1-\ell} \frac{(2(i-\ell)+j+1+|z|)^{2(i-\ell)+j+1}}{(i-\ell)!} e^{-\tilde{\beta} R e(z)} \\
\leq & C \tilde{\gamma}^{i+j+1} \frac{(2 i+j+1+|z|)^{2 i+j}}{(i-1)!} e^{-\tilde{\beta} R e(z)}\left\{\frac{1}{\tilde{\gamma}^2 i(2 i+j+1+|z|)}+\sum_{\ell=1}^{\infty}\left[\frac{\gamma_{\tilde{b}}^{\ell}}{\tilde{\gamma}^{\ell}}+\frac{\gamma_{\tilde{c}}^{\ell}}{\tilde{\gamma}^{\ell}}\right]\right\} .
\end{aligned}
$$
%
The assumptions on $\tilde{\gamma}$ allow us to infer that the geometric series converges to a bounded quantity, hence the expression in braces is also bounded and we have the necessary setup to use Lemma \ref{lemma_aux0}, which yields
%
\begin{eqnarray*}
\vert \tilde{u}_{i,j+1}(z) \vert &\leq& C e^{-\tilde{\beta} Re(z)} \left[ \tilde{\gamma}^{i+j+1} \frac{(2i+j+1+|z|)^{2i+j+1}}{(2i+j+1) (i-1)!}
+|{u}_{i,j+1}(0)|+|{u}'_{i,j+1}(0)| \right] \\
&\leq& C e^{-\tilde{\beta} Re(z)} \tilde{\gamma}^{i+j+1} \frac{(2i+j+1+|z|)^{2i+j+1}}{i!} \left[ \frac{i}{(2i+j+1)}+\frac{\tilde{C} i^{2i}}{\tilde{\gamma}^{j+1} (2i+j+1+|z|)^{2i+j+1}} \right] \\
&\leq& C e^{-\tilde{\beta} Re(z)} \tilde{\gamma}^{i+j+1} \frac{(2i+j+1+|z|)^{2i+j+1}}{i!}.
\end{eqnarray*}
%
Following the same steps as above, we can show the desired results for $\tilde{u}^{BL}_{i,0},\tilde{u}^{BL}_{i,1}, i\geq 1$. Now, we assume that the result holds for $\bar{u}^{BL}_{i,j}, j > i+1$ and we will establish it for $\bar{u}^{BL}_{i,j+1}$, with $i\geq 1$ fixed. By (\ref{ubar_ij}), the function $\bar{u}^{BL}_{i,j+1}$ satisfies
%
$$
\left(\bar{u}_{i, j+1}^{B L}\right)^{(4)}-\bar{b}_0\left(\bar{u}_{i, j+1}^{B L}\right)^{\prime \prime}=-\bar{c}_0 \bar{u}_{i, j-1}^{B L}+\sum_{\ell=1}^i\left\{\bar{b}_{\ell}\left(\bar{u}_{i-\ell, j+1-\ell}^{B L}\right)^{\prime \prime}+\bar{b}_{\ell}^{\prime}\left(\bar{u}_{i-\ell, j+1-\ell}^{B L}\right)^{\prime}-\bar{c}_{\ell} \bar{u}_{i-\ell, j-1-\ell}^{B L}\right\} .
$$
%
We bound the right hand side (which we call $\bar{G}$)  as follows, with the aid of the induction hypothesis and Cauchy's Integral Theorem for Derivatives:
$$
\begin{aligned}
& |\bar{G}(z)| \leq\left|\bar{c}_0\right|\left|\bar{u}_{i, j-1}^{B L}\right|+C \sum_{\ell=1}^i\left\{\left|\bar{b}_{\ell}\right|\left|\left(\bar{u}_{i-\ell, j+1-\ell}^{B L}\right)^{\prime \prime}\right|+\left|\bar{b}_{\ell}^{\prime}\right|\left|\left(\bar{u}_{i-\ell, j+1-\ell}^{B L}\right)^{\prime}\right|+\left|\bar{c}_{\ell}\right|\left|\bar{u}_{i-\ell, j-1-\ell}^{B L}\right|\right\} \\
& \leq C \bar{\gamma}^{i+j-1} \frac{(i+2(j-1)+|z|)^{i+2(j-1)}}{(j-1)!} e^{-\bar{\beta} R e(z)}+ \\
& +C \sum_{\ell=1}^i|z|^{\ell}\left(\gamma_{\bar{\ell}}^{\ell}\left|\bar{u}_{i-\ell, j+1-\ell}^{B L}\right|+\gamma_{\bar{c}}^{\ell}\left|\bar{u}_{i-\ell, j-1-\ell}^{B L}\right|\right) \\
& \leq C \bar{\gamma}^{i+j-1} \frac{(i+2(j-1)+|z|)^{i+2(j-1)}}{(j-1)!} e^{-\bar{\beta} R e(z)}+ \\
& +C \sum_{\ell=1}^i|z|^{\ell} \gamma_{\bar{b}}^{\ell} \bar{\gamma}^{i+j+1-2 \ell} \frac{(i-\ell+2(j+1-\ell)+|z|)^{i-\ell+2(j+1-\ell)}}{(j+1-\ell)!} e^{-\bar{\beta} R e(z)}+ \\
& +C \sum_{\ell=1}^i|z|^{\ell} \gamma_{\bar{c}}^{\ell} \bar{\gamma}^{i+j-1-2 \ell} \frac{(i-\ell+2(j-1-\ell)+|z|)^{i-\ell+2(j-1-\ell)}}{(j-1-\ell)!} e^{-\bar{\beta} R e(z)} \\
& \leq C \bar{\gamma}^{i+j+1} \frac{(i+2(j+1)+|z|)^{i+2 j+1}}{j!} e^{-\bar{\beta} R e(z)}\left\{\frac{j}{\bar{\gamma}^2(i+2(j+1)+|z|)^3}+\right. \\
& \left.+\frac{1}{\bar{\gamma}(i+2(j+1)+|z|)} \sum_{\ell=1}^{\infty} \frac{\gamma_{\bar{b}}^{\ell}}{\bar{\gamma}^{\ell}}+\frac{j(j-1)}{\bar{\gamma}^3(i+2(j+1)+|z|)^5} \sum_{\ell=1}^{\infty} \frac{\gamma_{\bar{c}}^{\ell}}{\bar{\gamma}^{\ell}}\right\} . \\
&
\end{aligned}
$$
%
The expression in braces is bounded, since both sums are convergent geometric series due to the assumptions on $\bar{\gamma}$, and an application of Lemma \ref{lemma_aux0} gives
%
$$
\begin{aligned}
\left|\bar{u}_{i, j+1}^{B L}(z)\right| \leq & C e^{-\bar{\beta} R e(z)}\left[\bar{\gamma}^{i+j+1} \frac{(i+2(j+1)+|z|)^{i+2(j+1)}}{(i+2(j+1)) j!}+\left|u_{i, j+1}(0)\right|+\left|u_{i, j+1}^{\prime}(0)\right|\right] \\
\leq & C e^{-\bar{\beta} R e(z)} \bar{\gamma}^{i+j+1} \frac{(i+2(j+1)+|z|)^{i+2(j+1)}}{(j+1)!} \times \\
& \times\left[\frac{j+1}{(i+2(j+1))}+\frac{\bar{C} i^i(j+1)^{j+1}}{\bar{\gamma}^{j+1}(i+2(j+1)+|z|)^{i+2(j+1)}}\right] \\
\leq & C e^{-\bar{\beta} R e(z)} \bar{\gamma}^{i+j+1} \frac{(i+2(j+1)+|z|)^{i+2(j+1)}}{(j+1)!} .
\end{aligned}
$$
%
Similarly, we can show the result for the remaining functions $\bar{u}_{i,j},i \geq 1$ and the proof is complete.
\end{proof}

\end{document}
